\documentclass[12pt,a4paper]{report}

\usepackage[left=3.0cm,right=2.25cm,top=2.5cm,bottom=2.4cm,headheight=15pt]{geometry}
\usepackage{mathptmx}
\usepackage[T1]{fontenc}
\usepackage{amsmath,amssymb}
\usepackage{setspace}
\usepackage{graphicx}
\usepackage{booktabs,array}
\usepackage{enumitem}
\usepackage{fancyhdr}
\usepackage{titlesec}
\usepackage{tikz}
\usetikzlibrary{arrows.meta,positioning,shapes.geometric,calc}
\usepackage[hidelinks]{hyperref}
\usepackage{caption}
\usepackage{tocloft}

% =====================================================================
%  GCOERC NASHIK - T.E. COMPUTER ENGINEERING SEMINAR REPORT TEMPLATE
%  Edit ONLY the 6 fields below, then write your chapters.
%  Keep logo.jpg (college) and sppu.jpg (university) in the same folder.
%  Compile with:  pdflatex report.tex  (run twice)  or  latexmk -pdf report.tex
% =====================================================================
\newcommand{\ttitle}{SEMINAR TITLE}            % title in normal case (used in header)
\newcommand{\TTITLE}{SEMINAR TITLE IN CAPITALS} % same title in CAPITALS (cover/certificates)
\newcommand{\studentname}{Student Name}
\newcommand{\rollno}{000}
\newcommand{\guidename}{Prof. Guide Name}
\newcommand{\acadyear}{2026--27}

\captionsetup{font=small,labelfont=bf}
\setlength{\parindent}{1.2em}
\setlength{\parskip}{0pt}

% ---------- chapter / section styling (as in college template) ----------
\titleformat{\chapter}[display]
  {\normalfont\bfseries\centering}{\large Chapter \thechapter}{10pt}{\large\uppercase}
\titlespacing*{\chapter}{0pt}{0pt}{16pt}
\titleformat{\section}{\normalfont\large\bfseries}{\thesection}{0.8em}{}
\titleformat{\subsection}{\normalfont\normalsize\bfseries}{\thesubsection}{0.8em}{}
\titlespacing*{\section}{0pt}{10pt}{4pt}
\titlespacing*{\subsection}{0pt}{8pt}{3pt}

% ---------- header / footer ----------

\fancypagestyle{main}{%
  \fancyhf{}
  \fancyhead[C]{\scriptsize \ttitle}
  \fancyfoot[C]{\small\thepage\\[1pt]\scriptsize Department of Computer Engineering, GCOERC, Nashik 2026--27}
  \renewcommand{\headrulewidth}{0.4pt}
  \renewcommand{\footrulewidth}{0pt}}
\pagestyle{main}
\fancypagestyle{plain}{\fancyhf{}%
  \fancyhead[C]{\scriptsize \ttitle}
  \fancyfoot[C]{\small\thepage\\[1pt]\scriptsize Department of Computer Engineering, GCOERC, Nashik 2026--27}
  \renewcommand{\headrulewidth}{0.4pt}}

\newcommand{\unnumchapter}[1]{%
  \chapter*{#1}\addcontentsline{toc}{chapter}{#1}}

% ---------- page border for cover and certificates ----------
\newcommand{\pageborder}{%
\begin{tikzpicture}[remember picture,overlay]
\draw[line width=1.6pt] ($(current page.north west)+(1.5cm,-1.5cm)$) rectangle ($(current page.south east)+(-1.5cm,1.5cm)$);
\end{tikzpicture}}

% ---------- TOC look: chapters bold, sections indented ----------
\renewcommand{\cftchapfont}{\bfseries\small}
\renewcommand{\cftchappagefont}{\bfseries\small}
\renewcommand{\cftsecfont}{\small}
\renewcommand{\cftsecpagefont}{\small}
\renewcommand{\cftsubsecfont}{\small}
\renewcommand{\cftsubsecpagefont}{\small}
\setlength{\cftbeforechapskip}{6pt}
\setlength{\cftbeforesecskip}{1pt}
\setcounter{tocdepth}{2}
\renewcommand{\cfttoctitlefont}{\hfill\normalfont\large\bfseries}
\renewcommand{\cftaftertoctitle}{\hfill}
\renewcommand{\cftloftitlefont}{\hfill\normalfont\large\bfseries}
\renewcommand{\cftafterloftitle}{\hfill}
\setlength{\cftbeforetoctitleskip}{0pt}
\setlength{\cftbeforeloftitleskip}{0pt}
\renewcommand{\cftchapdotsep}{\cftdotsep}
\renewcommand{\cftfigfont}{\small}
\renewcommand{\cftfigpagefont}{\small}

\begin{document}

% ------------------------------ COVER ------------------------------
\begin{titlepage}
\thispagestyle{empty}
\pageborder
\centering
\singlespacing
\vspace*{0.3cm}
{\large\bfseries A\par}\vspace{6pt}
{\large\bfseries SEMINAR REPORT\par}\vspace{6pt}
{\large\bfseries ON\par}\vspace{10pt}
{\Large\bfseries \TTITLE\par}
\vspace{14pt}
{\normalsize SUBMITTED TO THE SAVITRIBAI PHULE PUNE UNIVERSITY IN FULFILMENT OF THE REQUIREMENTS\\ FOR THE AWARD OF THE\par}
\vspace{8pt}
{\large\bfseries THIRD YEAR COMPUTER ENGINEERING\par}\vspace{8pt}
{\large\bfseries BY\par}
{\large\bfseries \studentname\par}
{\normalsize (Roll No: \rollno)\par}\vspace{8pt}
{\large\bfseries UNDER THE GUIDANCE OF\par}
{\large\bfseries \guidename\par}\vspace{10pt}
\includegraphics[height=3.6cm]{logo.jpg}\par\vspace{10pt}
{\normalsize\bfseries DEPARTMENT OF COMPUTER ENGINEERING\\
Guru Gobind Singh College of Engineering and Research Center, Nashik.\\
Khalsa Educational Complex, Guru Gobind Singh Marg, Wadala--Pathardi Road, Indira Nagar, Nashik, Maharashtra 422009\\
YEAR \acadyear\par}
\end{titlepage}

% --------------------------- CERTIFICATE 1 -------------------------
\thispagestyle{empty}
\pageborder
\begin{center}
\singlespacing
\vspace*{0.4cm}
{\normalsize\bfseries DEPARTMENT OF COMPUTER ENGINEERING\\
Guru Gobind Singh College of Engineering and Research Center, Nashik\\
Khalsa Educational Complex, Guru Gobind Singh Marg, Wadala--Pathardi Road, Indira Nagar, Nashik, Maharashtra 422009\\
Year \acadyear\par}
\vspace{10pt}
\includegraphics[height=3.4cm]{logo.jpg}\\[14pt]
{\large\bfseries CERTIFICATE}\\[8pt]
{\bfseries This is to certify that seminar report entitled}\\[10pt]
{\large\bfseries ``\TTITLE''}\\[10pt]
is submitted as fulfilment of curriculum of the T.E.\ Computer Engineering\\[10pt]
{\bfseries BY}\\[8pt]
{\bfseries \studentname}\\
{\bfseries (Roll No: \rollno)}
\end{center}
\vspace{40pt}
\begin{center}
\singlespacing\small
\begin{tabular}{p{0.28\textwidth}p{0.36\textwidth}p{0.24\textwidth}}
\guidename & Dr.\ Swati Khokale \& Mr.\ Ajit Pagar & Dr.\ Nita Thakare\\
\textbf{Seminar Guide} & \textbf{Seminar Coordinator} & \textbf{Head of Dept.}
\end{tabular}
\end{center}
\vfill
{\singlespacing\noindent\hspace{0.6cm}\textbf{Place:} GCOERC, Nashik\\
\hspace*{0.6cm}\textbf{Date:}\par}
\vspace{1.4cm}
\clearpage

% --------------------------- CERTIFICATE 2 -------------------------
\thispagestyle{empty}
\pageborder
\begin{center}
\singlespacing
\vspace*{0.4cm}
{\normalsize Savitribai Phule Pune University\par}\vspace{6pt}
\includegraphics[height=3.4cm]{sppu.jpg}\\[12pt]
{\large\bfseries CERTIFICATE}\\[10pt]
This is to certify that\\[8pt]
{\bfseries \studentname}\\
{\bfseries (Roll No: \rollno)}\\[10pt]
Student of T.E.\ Computer\\[4pt]
Was examined in Seminar Report entitled\\[10pt]
{\large\bfseries ``\TTITLE''}\\[12pt]
on \underline{\hspace{0.8cm}} / \underline{\hspace{0.8cm}} /2026\\[8pt]
At\\[4pt]
DEPARTMENT OF COMPUTER ENGINEERING,\\
GURU GOBIND SINGH COLLEGE OF ENGINEERING AND RESEARCH CENTER, NASHIK\\[4pt]
{\bfseries YEAR \acadyear}
\end{center}
\vspace{50pt}
\begin{center}
\begin{tabular}{p{0.38\textwidth}p{0.38\textwidth}}
\dotfill & \dotfill\\
\centering Internal Examiner & \centering External Examiner\tabularnewline
\end{tabular}
\end{center}
\clearpage

\pagenumbering{roman}
\setcounter{page}{1}
\onehalfspacing

% ------------------------------ ABSTRACT ---------------------------
\unnumchapter{ABSTRACT}
Write a 150--250 word abstract here: the problem, the approach studied in the seminar, the main findings reported in the source paper(s), and the conclusion. Cite sources with \verb|\cite{key}|.

\vspace{6pt}
\noindent\textbf{\emph{Keywords---keyword 1; keyword 2; keyword 3; keyword 4.}}
\clearpage

% -------------------------- ACKNOWLEDGEMENT ------------------------
\unnumchapter{ACKNOWLEDGEMENT}
It is my immense pleasure to work on this seminar \textbf{\ttitle}. It is only the blessing of my divine master which has prompted and mentally equipped me to undergo the study of this seminar.

I would like to thank \textbf{Prof.\ (Dr.) N.\ G.\ Nikam}, Principal, Guru Gobind Singh College of Engineering and Research Center, for giving me such an opportunity to develop practical knowledge about the subject. I am also thankful to \textbf{Dr.\ Nita M.\ Thakare}, Head of Computer Engineering Department, for her valuable encouragement at every phase of my seminar work and completion.

I offer my sincere thanks to my guide \textbf{\guidename}, who very affectionately encouraged me to work on the subject and gave me valuable guidance from time to time. While preparing this seminar I am very much thankful to her/him.

I am also grateful to the entire staff of Computer Engineering Department for their kind co-operation, which helped me in the successful completion of this seminar.

\vspace{30pt}
\hfill\begin{minipage}{0.45\textwidth}
\singlespacing
\textbf{\studentname}\\
\textbf{Roll No: \rollno}\\[6pt]
GCOERC, NASHIK.
\end{minipage}
\clearpage

% ------------------------------- INDEX -----------------------------
\begingroup
\singlespacing
\renewcommand{\contentsname}{INDEX}
\addcontentsline{toc}{chapter}{INDEX}
\tableofcontents
\clearpage
\renewcommand{\listfigurename}{LIST OF FIGURES}
\addcontentsline{toc}{chapter}{LIST OF FIGURES}
\listoffigures
\endgroup
\clearpage

\pagenumbering{arabic}
\setcounter{page}{1}
\onehalfspacing

% ----------------------------- CHAPTER 1 ---------------------------
\chapter{INTRODUCTION}
Introduce the topic: what it is, why it matters, and the background needed to follow the report \cite{ref1}.

\begin{figure}[htbp]
\centering
\fbox{\parbox[c][4cm][c]{0.8\textwidth}{\centering Insert your figure here using\\ \texttt{\textbackslash includegraphics[width=0.8\textbackslash textwidth]\{figure1.png\}}}}
\caption{Caption of the figure.}\label{fig:one}
\end{figure}

\section{Motivation}
Why was this topic chosen? What real-world problem motivates it?

\section{Problem Statement}
State the problem addressed in one clear paragraph.

% ----------------------------- CHAPTER 2 ---------------------------
\chapter{OBJECTIVES}
The main objectives of this seminar are:
\begin{enumerate}[leftmargin=*]
  \item To understand the fundamental concepts of \ldots
  \item To study the working methodology and architecture of \ldots
  \item To explore its applications in real-world scenarios.
  \item To analyse its advantages, limitations and challenges.
  \item To study future developments and scope in the selected field.
\end{enumerate}

% ----------------------------- CHAPTER 3 ---------------------------
\chapter{LITERATURE REVIEW}
\noindent\textbf{Title of Paper 1}\\
Author et al.\ \cite{ref1} proposed \ldots\ Key findings \ldots\ Limitation \ldots

\noindent\textbf{Title of Paper 2}\\
Author et al.\ \cite{ref2} studied \ldots

\noindent\textbf{Title of Paper 3}\\
Author et al.\ \cite{ref3} presented \ldots

% ----------------------------- CHAPTER 4 ---------------------------
\chapter{METHODOLOGY / SYSTEM ARCHITECTURE}
Explain the working of the system step by step.

\section{System Architecture}
Describe the architecture and refer to the diagram (Figure~\ref{fig:arch}).

\begin{figure}[htbp]
\centering
\begin{tikzpicture}[node distance=1.4cm,
  box/.style={draw,rounded corners,minimum width=2.8cm,minimum height=1cm,align=center,font=\small},
  arr/.style={-{Stealth[length=2.5mm]},thick}]
\node[box,fill=blue!10] (a) {Input};
\node[box,fill=green!10,right=of a] (b) {Processing};
\node[box,fill=orange!15,right=of b] (c) {Output};
\draw[arr] (a)--(b); \draw[arr] (b)--(c);
\end{tikzpicture}
\caption{Block diagram of the system.}\label{fig:arch}
\end{figure}

\section{Working}
\begin{enumerate}[leftmargin=*]
  \item Step one \ldots
  \item Step two \ldots
  \item Step three \ldots
\end{enumerate}

\begin{table}[htbp]
\centering\small\singlespacing
\caption{Sample table}\label{tab:one}
\begin{tabular}{|p{4cm}|p{8cm}|}
\hline
\textbf{Parameter} & \textbf{Value / Description}\\\hline
Item 1 & Description\\\hline
Item 2 & Description\\\hline
\end{tabular}
\end{table}

An equation example:
\begin{equation}
y = f(x) + \varepsilon .
\end{equation}

% ----------------------------- CHAPTER 5 ---------------------------
\chapter{APPLICATIONS}
The technology discussed in this seminar has applications in various domains.

\section{Application 1 -- Name of Domain}
Explain how the technology is used, the problem it addresses and the benefits.

\section{Application 2 -- Name of Domain}
Describe its role, functionality and practical significance.

\section{Application 3 -- Name of Domain}
Explain how it can be implemented and what advantages it provides.

\section{Application 4 -- Name of Domain}
Discuss another important real-world application.

% ----------------------------- CHAPTER 6 ---------------------------
\chapter{ADVANTAGES AND LIMITATIONS}
\section{Advantages}
\begin{itemize}[leftmargin=*]
  \item Advantage 1.
  \item Advantage 2.
  \item Advantage 3.
\end{itemize}
\section{Limitations}
\begin{itemize}[leftmargin=*]
  \item Limitation 1.
  \item Limitation 2.
  \item Limitation 3.
\end{itemize}

% ----------------------------- CHAPTER 7 ---------------------------
\chapter{CONCLUSION}
Summarise what was studied, the main findings, the limits of the work, and the future scope.

% ----------------------------- REFERENCES --------------------------
\renewcommand{\bibname}{References}
\begin{singlespacing}
\begin{thebibliography}{99}
\addcontentsline{toc}{chapter}{REFERENCES}
\bibitem{ref1} A. Author, B. Author, and C. Author, ``Title of the paper,'' \emph{Journal Name}, vol. X, no. Y, pp. 1--10, Year, doi: 10.xxxx/xxxxx.
\bibitem{ref2} A. Author \emph{et al.}, ``Title of the paper,'' in \emph{Proc. Conference Name}, Year, pp. 1--8.
\bibitem{ref3} A. Author, ``Title of the paper,'' arXiv:xxxx.xxxxx, Year.
\end{thebibliography}
\end{singlespacing}

\end{document}
